\documentclass[11pt,letterpaper]{article}

\usepackage[utf8]{inputenc}
\usepackage[margin=1in]{geometry}
\usepackage{enumitem}
\usepackage{mathpazo}
\usepackage[T1]{fontenc}
\usepackage{xcolor}
\usepackage{fancyhdr}

\pagestyle{fancy}
\fancyhf{}
\fancyhead[L]{\footnotesize\textcolor{gray!90}{\textit{Synergistic Activities}}}
\fancyhead[R]{\footnotesize\textcolor{gray!90}{\textit{Mostafa Rezaali}}}
\renewcommand{\headrulewidth}{0.4pt}
\renewcommand{\footrulewidth}{0pt}

\setlength{\parindent}{0pt}
\setlength{\parskip}{4pt}
\setlist[enumerate]{leftmargin=*, itemsep=5pt, topsep=4pt, parsep=0pt}

\begin{document}

\begin{center}
    {\Large\bfseries Synergistic Activities}\\[0.12cm]
    {\normalsize Mostafa Rezaali}\\[-0.02cm]
    {\small Ph.D. Candidate, Climate Science (Geography), University of Florida}\\[0.08cm]
    \rule{0.62\textwidth}{0.4pt}
\end{center}

\begin{enumerate}
\item \textbf{Mentoring emerging climate-data scientists.}
Served as an NSF LEAP REU mentor in 2025 and supported undergraduate and master's-level research using neural-network models for climate and environmental hazards. This work helped students turn climate-risk questions into reproducible data workflows, model experiments, uncertainty summaries, and interpretable results.

\item \textbf{Teaching deep learning for climate science.}
Delivered invited lectures at the University of Florida on deep learning applications in climate science, including CNN, LSTM, and hybrid ConvLSTM approaches for spatiotemporal climate data. These lectures connected machine-learning concepts to climate impacts, geospatial analysis, and responsible use of artificial intelligence for environmental prediction.

\item \textbf{Transferring AI methods across climate, hydrology, and air quality.}
Developed and published machine-learning workflows for tropospheric ozone forecasting, urban water-demand forecasting, air-dispersion modeling, and environmental-health risk assessment. This cross-domain work supports broader transfer of atmospheric-science methods into decision-relevant environmental data science.

\item \textbf{Professional peer review and disciplinary service.}
Reviewed manuscripts for interdisciplinary journals including \textit{Scientific Reports}, \textit{Journal of Hydrology}, and \textit{Springer Nature Applied Sciences}. This service contributes to rigorous scholarship at the intersection of climate science, hydrology, geospatial analysis, environmental modeling, and applied machine learning.

\item \textbf{Open and reproducible climate-risk workflows.}
Built Python, MATLAB, R, and geospatial workflows for large environmental datasets, including NetCDF/GRIB processing, model training, hindcast evaluation, and diagnostic visualization. The proposed fellowship will extend this commitment through transparent evaluation of probabilistic forecasts and reusable workflows for heat-extreme prediction.
\end{enumerate}

\end{document}
